\scanpage{362}\mspage{278}
\uline{Corollaire} \deleted{\unsure{3|2}}4.12.4. Soient $M$ et $N$ deux
catégories accessibles, et soit $\pi_0$ un cardinal
infini tel que $M$ soit $\pi_0$-accessible, i.e.
que $\pi_0$ soit un cardinal utile pour $M$.
Soit $\pi_1=\operatorname{Sup}(\pi_0,\operatorname{Card\,red}M_{\pi_0})$. Alors
pour tt cardinal \inserted{$\pi\geq\pi_1$} utile pour $M$
(et en particulier, pour tt $\pi\geq\pi_1$ tel
que $\pi^{\pi_0}=\pi$), la catégorie des foncteurs
$M\unsure{-|\longrightarrow}N$ qui sont $\pi$-accessibles est
elle-\unsure{$=$} accessible.

En effet, on aura $M\approx\mathrm{Ind}_\pi M_\pi$,
d'où par 4.12.3
\[
 \underline{\mathrm{Hom}}_{!\pi}(M,N)\xrightarrow{\approx}\underline{\mathrm{Hom}}(M_\pi,N),
\]
or $\underline{\mathrm{Hom}}(M_\pi,N)$ est accessible par
4.9.5 (\unsure{p|p\textsuperscript{e}|p\textsuperscript{ge}} 220). On peut préciser
d'ailleurs ce point. Supposons $\pi_0$ choisi
tel que $\pi_0$ soit utile non seulemt pour
$M$, mais pour $N$ égalemt, \deleted{et supposons de
\unsure{$=$} $\pi_1\geq\operatorname{Sup}(\operatorname{Card\,red}M_{\pi_0},\operatorname{Card\,red}N_{\pi_0})$.}
\deleted{\unsure{Notons}}\inserted{Rappelons}
que l'on a $\operatorname{Card\,Fl} M_\pi\leq 2^\pi$ (cf. \unsure{4.2.7.1|4.2.7.})
page 173)
% archive sheet number: 361

